% Topic 1.2.09 — Markov Processes and Discrete Markov Chains.

\section{Markov Processes and Discrete Markov Chains}
\label{sec:1.2.09-markov-chains}

\begin{ticket}
Markov stochastic process. Discrete Markov chain. Transition probabilities; Kolmogorov–Chapman equations. Homogeneous discrete Markov chains.
\end{ticket}

\subsection*{Theory}

A Markov process is a stochastic process in which the future depends
on the past only through the present. This is the simplest, and by
far the most common, way to build dependent processes from a recipe
for ``one step'' transitions.

\begin{definition}[Markov property]
\label{def:markov-property}
A discrete-time stochastic process $(X_n)_{n \geq 0}$ taking values
in a countable state space~$S$ has the \emph{Markov property} if for
every $n \geq 0$, every $i_0, \dots, i_{n+1} \in S$, and whenever
the conditioning event has positive probability,
\[
  \Prob(X_{n+1} = i_{n+1} \mid X_n = i_n, \dots, X_0 = i_0) \;=\; \Prob(X_{n+1} = i_{n+1} \mid X_n = i_n).
\]
A process satisfying this property is a \emph{(discrete-time) Markov
chain} on~$S$.
\end{definition}

\begin{definition}[Transition probabilities; homogeneity]
\label{def:transition-probabilities}
The \emph{transition probability} from state~$i$ to state~$j$ at
step~$n$ is
\[
  p_{ij}^{(n,n+1)} = \Prob(X_{n+1} = j \mid X_n = i).
\]
The chain is \emph{homogeneous} (or \emph{time-homogeneous}) if
$p_{ij}^{(n,n+1)} = p_{ij}$ does not depend on~$n$. The
\emph{transition matrix} is $\mat{P} = (p_{ij})_{i, j \in S}$; it is
\emph{stochastic}, meaning $p_{ij} \geq 0$ and $\sum_j p_{ij} = 1$
for every~$i$.
\end{definition}

By the multiplication rule (\cref{prop:multiplication-rule}), the
joint law of $(X_0, X_1, \dots, X_n)$ on a homogeneous chain
factorises as
\[
  \Prob(X_0 = i_0, X_1 = i_1, \dots, X_n = i_n) \;=\; \Prob(X_0 = i_0) \cdot p_{i_0 i_1} \cdot p_{i_1 i_2} \cdots p_{i_{n-1} i_n},
\]
so the initial distribution $\pi_0$ on~$S$ together with the
transition matrix $\mat{P}$ determines all FDDs of the chain
(\cref{def:fdd}); via Kolmogorov's existence theorem
(\cref{thm:kolmogorov-extension}) this is enough to construct the
chain on the path space~$S^{\N}$.

\begin{definition}[$n$-step transition probabilities]
\label{def:n-step-transition}
For a homogeneous Markov chain, the \emph{$n$-step transition
probability} is
\[
  p_{ij}^{(n)} \;=\; \Prob(X_n = j \mid X_0 = i),
\]
with the convention $p_{ij}^{(0)} = \delta_{ij}$ (Kronecker delta).
The matrix $\mat{P}^{(n)} = (p_{ij}^{(n)})_{i,j \in S}$ is also
stochastic.
\end{definition}

\begin{theorem}[Kolmogorov–Chapman equations]
\label{thm:kolmogorov-chapman}
For a homogeneous Markov chain and any $m, n \geq 0$,
\[
  p_{ij}^{(n + m)} \;=\; \sum_{k \in S} p_{ik}^{(n)}\, p_{kj}^{(m)} \qquad \text{for all } i, j \in S,
\]
or equivalently in matrix form,
$\mat{P}^{(n+m)} = \mat{P}^{(n)}\, \mat{P}^{(m)}$. In particular,
$\mat{P}^{(n)} = \mat{P}^n$ — the $n$-step transition matrix is the
$n$-th matrix power of the one-step transition matrix.
\end{theorem}
\proofof{kolmogorov-chapman}

\begin{corollary}[Distribution at time~$n$]
\label{cor:distribution-at-n}
If~$X_0$ has distribution $\pi_0 = (\pi_0(i))_{i \in S}$ (a row
vector), then $X_n$ has distribution $\pi_n = \pi_0\, \mat{P}^n$.
\end{corollary}
\proofof{distribution-at-n}

\medskip
\noindent\textbf{Stationary distribution.} A probability vector
$\pi$ on $S$ is \emph{stationary} (or \emph{invariant}) if
$\pi = \pi\, \mat{P}$, i.e.\ $\pi(j) = \sum_i \pi(i) p_{ij}$ for every
$j \in S$. By induction, $\pi = \pi\, \mat{P}^n$ for all
$n \geq 0$, so a chain started in a stationary distribution remains
in it for all time. For finite-state chains, existence of a
stationary distribution follows from the Perron–Frobenius theorem
applied to $\mat{P}^{\T}$; uniqueness and convergence
$\pi_0 \mat{P}^n \to \pi$ require irreducibility and aperiodicity,
the standard ergodic theorem for Markov chains
\cite[Ch.~VIII]{Shiryaev1989}.

\begin{remark}
Continuous-time Markov processes — for which the transition
mechanism is an infinitesimal generator rather than a transition
matrix — are constructed analogously. The competing-exponentials
problem (\cref{prob:exponential-min}) is the elementary case: the
holding time in each state is exponential, and at the time of a
transition the next state is selected by the chain on jumps. See
\cite[Ch.~VIII]{Sevastyanov1982}.
\end{remark}

\subsection*{Proofs}

\begin{proof}[\Cref{thm:kolmogorov-chapman}]
\proofanchor{kolmogorov-chapman}
By the law of total probability (\cref{thm:total-probability})
applied to the partition $\{X_n = k\}_{k \in S}$,
\[
  \Prob(X_{n+m} = j \mid X_0 = i) = \sum_{k \in S} \Prob(X_{n+m} = j,\, X_n = k \mid X_0 = i).
\]
By the multiplication rule, each term equals
$\Prob(X_n = k \mid X_0 = i) \cdot \Prob(X_{n+m} = j \mid X_n = k, X_0 = i)$.
The Markov property collapses the second factor:
$\Prob(X_{n+m} = j \mid X_n = k, X_0 = i) = \Prob(X_{n+m} = j \mid X_n = k) = p_{kj}^{(m)}$
(the rightmost equality uses time-homogeneity, since the conditional
distribution of the chain over $m$ further steps starting from
state~$k$ at any time is the same). Combining,
\[
  p_{ij}^{(n+m)} = \sum_{k} p_{ik}^{(n)} p_{kj}^{(m)}, \qquad \mat{P}^{(n+m)} = \mat{P}^{(n)} \mat{P}^{(m)}.
\]
Setting $\mat{P}^{(0)} = \mat{I}$ and iterating from $\mat{P}^{(1)} = \mat{P}$
gives $\mat{P}^{(n)} = \mat{P}^n$ by induction.
\end{proof}

\begin{proof}[\Cref{cor:distribution-at-n}]
\proofanchor{distribution-at-n}
By total probability conditioning on the initial state,
\[
  \pi_n(j) = \Prob(X_n = j) = \sum_{i \in S} \Prob(X_0 = i) \Prob(X_n = j \mid X_0 = i) = \sum_i \pi_0(i) p_{ij}^{(n)},
\]
which is the $j$-th component of the row-vector–matrix product
$\pi_0 \mat{P}^{(n)} = \pi_0 \mat{P}^n$.
\end{proof}

\subsection*{Examples}

\begin{example}[Two-state chain: weather]
\label{ex:two-state-chain}
Let $S = \{R, S\}$ (rainy, sunny) and
\[
  \mat{P} = \begin{pmatrix} 0.7 & 0.3 \\ 0.4 & 0.6 \end{pmatrix},
\]
where rows index the current state and columns the next. To find
the probability of rain in two days given that today is sunny,
read off $p_{SR}^{(2)} = (\mat{P}^2)_{SR}$:
\[
  \mat{P}^2 = \begin{pmatrix} 0.61 & 0.39 \\ 0.52 & 0.48 \end{pmatrix},
\]
giving $p_{SR}^{(2)} = 0.52$. The unique stationary distribution
solves $\pi = \pi \mat{P}$ with $\pi(R) + \pi(S) = 1$:
$0.3\, \pi(R) = 0.4\, \pi(S)$, so $\pi(R) = 4/7$ and $\pi(S) = 3/7$ —
the long-run frequency of rainy days.
\end{example}

\begin{example}[Simple random walk on $\Z$]
\label{ex:rw-on-Z}
$X_{n+1} = X_n + \xi_{n+1}$ with $\xi_n$ i.i.d.\ taking $+1$ with
probability~$p$ and $-1$ with probability~$1 - p$. The state space
is $\Z$, the chain is homogeneous, and the transition probabilities
are
\[
  p_{i, i+1} = p, \qquad p_{i, i-1} = 1 - p, \qquad p_{i, j} = 0 \text{ otherwise}.
\]
Two-step transitions: $p_{i, i}^{(2)} = 2 p (1 - p)$,
$p_{i, i+2}^{(2)} = p^2$, $p_{i, i-2}^{(2)} = (1-p)^2$. The
random walk visits any specific state with probability~$1$ when
$p = 1/2$ (recurrence) and with probability strictly less
than~$1$ when $p \neq 1/2$ (transience); the cleanest proof uses
generating-function arguments \cite[Ch.~XII]{Sevastyanov1982}.
\end{example}

\begin{example}[Ehrenfest urn]
\label{ex:ehrenfest}
$N$~balls are distributed between two urns. At each step, one ball
is chosen uniformly and moved to the other urn. Let $X_n$ be the
number of balls in the first urn. Then $S = \{0, 1, \dots, N\}$
and
\[
  p_{i, i-1} = \frac{i}{N}, \qquad p_{i, i+1} = \frac{N - i}{N}, \qquad p_{i, j} = 0 \text{ otherwise}.
\]
A direct calculation shows that the binomial distribution
$\pi(i) = \binom{N}{i} 2^{-N}$ is stationary. The Ehrenfest model
illustrates how a deterministic-looking deterministic-feel Boltzmann
``equilibrium'' arises from a mixing Markov chain on a finite
state space.
\end{example}

\begin{problem}
\label{prob:absorbing-chain}
A gambler with capital $i \in \{0, 1, \dots, N\}$ wins or loses
$1$~unit per round independently with probability $1/2$ each, until
hitting~$0$ (ruin) or~$N$ (jackpot). Compute the probability
$h_i = \Prob(\text{hit } N \text{ before } 0 \mid X_0 = i)$.
\end{problem}

\begin{proof}[Solution]
States $0$ and $N$ are absorbing ($p_{00} = p_{NN} = 1$). For the
interior states the Markov property gives the linear recurrence
$h_i = \tfrac{1}{2} h_{i+1} + \tfrac{1}{2} h_{i-1}$ for
$1 \leq i \leq N - 1$, with boundary $h_0 = 0$ and $h_N = 1$. The
recurrence says $h_{i+1} - h_i = h_i - h_{i-1}$, so the differences
are constant; denote them by~$c$. Then $h_N - h_0 = N c = 1$, so
$c = 1/N$ and
\[
  h_i \;=\; \frac{i}{N}.
\]
For asymmetric step probabilities $p \neq 1/2$, the same conditioning
gives a non-trivial geometric solution
$h_i = (1 - r^i)/(1 - r^N)$ with $r = (1 - p)/p$ — the classical
\emph{gambler's ruin} formula \cite[Ch.~VI, \S~3]{Shiryaev1989}.
\end{proof}
